
\documentclass{article}
\usepackage{colortbl}
\usepackage{makecell}
\usepackage{multirow}
\usepackage{supertabular}

\begin{document}

\newcounter{utterance}

\centering \large Interaction Transcript for game `hot\_air\_balloon', experiment `air\_balloon\_survival\_en\_reasoning off\_hard', episode 2 with Qwen3{-}32B{-}t0.0.
\vspace{24pt}

{ \footnotesize  \setcounter{utterance}{1}
\setlength{\tabcolsep}{0pt}
\begin{supertabular}{c@{$\;$}|p{.15\linewidth}@{}p{.15\linewidth}p{.15\linewidth}p{.15\linewidth}p{.15\linewidth}p{.15\linewidth}}
    \# & $\;$A & \multicolumn{4}{c}{Game Master} & $\;\:$B\\
    \hline

    \theutterance \stepcounter{utterance}  
    & & \multicolumn{4}{p{0.6\linewidth}}{
        \cellcolor[rgb]{0.9,0.9,0.9}{
            \makecell[{{p{\linewidth}}}]{
                \texttt{\tiny{[P1$\langle$GM]}}
                \texttt{You are participating in a collaborative negotiation game.} \\
\\ 
\texttt{Together with another participant, you must agree on a single set of items that will be kept. Each of you has your own view of how much each item matters to you (importance). You do not know how the other participant values the items. Additionally, you are given the effort each item demands.} \\
\texttt{You may only agree on a set if the total effort of the selected items does not exceed a shared limit:} \\
\\ 
\texttt{LIMIT = 3203} \\
\\ 
\texttt{Here are the individual item effort values:} \\
\\ 
\texttt{Item effort = \{"A21": 52, "B93": 577, "A11": 177, "B33": 702, "C08": 119, "B39": 232, "B96": 577, "C85": 205, "A82": 516, "B31": 582, "B20": 676, "A60": 906, "B09": 316, "C81": 433, "B21": 336\}} \\
\\ 
\texttt{Here is your personal view on the importance of each item:} \\
\\ 
\texttt{Item importance values = \{"A21": 138, "B93": 583, "A11": 868, "B33": 822, "C08": 783, "B39": 65, "B96": 262, "C85": 121, "A82": 508, "B31": 780, "B20": 461, "A60": 484, "B09": 668, "C81": 389, "B21": 808\}} \\
\\ 
\texttt{Goal:} \\
\\ 
\texttt{Your goal is to negotiate a shared set of items that benefits you as much as possible (i.e., maximizes total importance to YOU), while staying within the LIMIT. You are not required to make a PROPOSAL in every message {-} you can simply negotiate as well. All tactics are allowed!} \\
\\ 
\texttt{Interaction Protocol:} \\
\\ 
\texttt{You may only use the following structured formats in a message:} \\
\\ 
\texttt{PROPOSAL: \{'A', 'B', 'C', …\}} \\
\texttt{Propose keeping exactly those items.} \\
\texttt{REFUSE: \{'A', 'B', 'C', …\}} \\
\texttt{Explicitly reject opponent's proposal.} \\
\texttt{ARGUMENT: \{'...'\}} \\
\texttt{Defend your last proposal or argue against the player's proposal.} \\
\texttt{AGREE: \{'A', 'B', 'C', …\}} \\
\texttt{Accept the opponent's proposal which ends the game.} \\
\\ 
\\ 
\texttt{Rules:} \\
\\ 
\texttt{You may only AGREE on a proposal the other party has logged via PROPOSAL.} \\
\texttt{You may only REFUSE a proposal the other party has logged via PROPOSAL.} \\
\texttt{Total effort of any PROPOSAL or AGREE set must be ≤ LIMIT.} \\
\texttt{Do NOT reveal your hidden importance scores.} \\
\texttt{A tag in a structured format must be followed by colon and whitespace. The argument must be a python set containing 0 or more strings.} \\
\texttt{So, it must be of the form TAG: \{...\}} \\
\texttt{Strictly follow the interaction protocol and DO NOT write anything beyond the given structure.} \\
\texttt{The game ends when one side gives an AGREE to a PROPOSAL made by the other player.} \\
\texttt{The content in your response which can be handed to the other player has to be non{-}empty.} \\
\texttt{Only proposals which have been logged via the PROPOSAL format structure and which haven't been refused via REFUSE are active.} \\
\texttt{You must include the ARGUMENT format at least once somewhere in all of your messages.} \\
            }
        }
    }
    & \\ \\

    \theutterance \stepcounter{utterance}  
    & \multicolumn{4}{p{0.6\linewidth}}{
        \cellcolor[rgb]{0.9,0.9,0.9}{
            \makecell[{{p{\linewidth}}}]{
                \texttt{\tiny{[P1$\rangle$GM]}}
                \texttt{<think>} \\
\texttt{Okay, let's see. I need to negotiate a set of items that maximizes my total importance without exceeding the effort limit of 3203. First, I should figure out which items are most important to me and check their combined effort.} \\
\\ 
\texttt{Looking at my importance values: A11 is the highest at 868, followed by B33 (822), B21 (808), B31 (780), C08 (783), B09 (668), A60 (484), A82 (508), B20 (461), B93 (583), B96 (262), C81 (389), C85 (121), B39 (65), A21 (138). The lowest importance items are B39 and A21.} \\
\\ 
\texttt{Now, the effort for each item. Let's list them out. A11 is 177, B33 is 702, B21 is 336, B31 is 582, C08 is 119, B09 is 316, A60 is 906, A82 is 516, B20 is 676, B93 is 577, B96 is 577, C81 is 433, C85 is 205, B39 is 232, A21 is 52.} \\
\\ 
\texttt{I need to select a combination where the total effort is ≤3203. Let's start by including the highest importance items and see if we can fit them.} \\
\\ 
\texttt{Start with A11 (177), B33 (702), B21 (336), B31 (582), C08 (119), B09 (316). Let's add these up: 177+702=879, +336=1215, +582=1797, +119=1916, +316=2232. That's 2232. Remaining effort:} \\
            }
        }
    }
    & & \\ \\

    \theutterance \stepcounter{utterance}  
    & & \multicolumn{4}{p{0.6\linewidth}}{
        \cellcolor[rgb]{0.9,0.9,0.9}{
            \makecell[{{p{\linewidth}}}]{
                \texttt{\tiny{[P1$\langle$GM]}}
                \texttt{Your response did not contain an argument. You must include the structured format ARGUMENT: \{...\} somewhere in your response. Try again} \\
            }
        }
    }
    & \\ \\

    \theutterance \stepcounter{utterance}  
    & \multicolumn{4}{p{0.6\linewidth}}{
        \cellcolor[rgb]{0.9,0.9,0.9}{
            \makecell[{{p{\linewidth}}}]{
                \texttt{\tiny{[P1$\rangle$GM]}}
                \texttt{<think>} \\
\texttt{Okay, let me try again. The user mentioned that I need to include an ARGUMENT in my response. So, first, I need to make sure that I structure my message correctly with the required tags.} \\
\\ 
\texttt{Looking back at the problem, my goal is to propose a set of items that maximizes my total importance without exceeding the effort limit. I had started by selecting the highest importance items and calculating their total effort. The initial total was 2232, leaving 971 effort left. I can add more items within that remaining effort.} \\
\\ 
\texttt{Let me check which other high{-}importance items I can include. The next ones are A60 (484 importance) with effort 906. Wait, 906 is more than the remaining 971, but adding it would take the total to 2232 + 906 = 3138, which is under 3203. That's possible. Then, maybe add A82 (508 importance) with effort 516. Adding that would make total effort 3138 + 516 = 3654, which exceeds the limit. So that's too much. Alternatively, maybe C81 (389 importance) with effort 433. 2232 + 433 = 2665, which is still under. Then maybe B93 (583 importance) with effort 577. 2665 + 577 = 3242, which is over 3203. So that's too much. Hmm.} \\
\\ 
\texttt{Alternatively, after adding A60 (906 effort), total becomes 2232 + 906 = 3138. Then, check if there's any other item that can fit. The remaining is 3203 {-} 3138 = 65. The smallest effort items are A21 (52) and B39 (232). A21 can be added. So total effort would be 3138 +52=3190, which is under. Then, maybe add C85 (121 importance) with effort 205. But 3190 +205=3395, which} \\
            }
        }
    }
    & & \\ \\

    \theutterance \stepcounter{utterance}  
    & & \multicolumn{4}{p{0.6\linewidth}}{
        \cellcolor[rgb]{0.9,0.9,0.9}{
            \makecell[{{p{\linewidth}}}]{
                \texttt{\tiny{[P1$\langle$GM]}}
                \texttt{Your response did not contain an argument. You must include the structured format ARGUMENT: \{...\} somewhere in your response. Try again} \\
            }
        }
    }
    & \\ \\

    \theutterance \stepcounter{utterance}  
    & \multicolumn{4}{p{0.6\linewidth}}{
        \cellcolor[rgb]{0.9,0.9,0.9}{
            \makecell[{{p{\linewidth}}}]{
                \texttt{\tiny{[P1$\rangle$GM]}}
                \texttt{<think>} \\
\texttt{Okay, I need to make sure I include an ARGUMENT in my response. Let me start by recalling the previous steps. The user wants me to propose a set of items that maximizes my total importance without exceeding the effort limit. I had started by selecting high{-}importance items and calculating their total effort. The initial set was A11, B33, B21, B31, C08, B09 with a total effort of 2232. Then I considered adding A60 (effort 906) which would bring the total to 3138, leaving 65 effort. I could add A21 (52) to reach 3190. Then maybe C85 (205) but that would exceed. Alternatively, maybe another item.} \\
\\ 
\texttt{Wait, the user's last message said my previous response didn't have an argument. So I need to structure my next message with PROPOSAL and ARGUMENT. Let me check the rules again. The ARGUMENT must be a Python set of strings, but the content is the actual argument. Wait, no, the ARGUMENT tag is followed by a set of strings? Or is it a set of items? Wait, the instructions say: "The argument must be a python set containing 0 or more strings." So the ARGUMENT is a set of strings, which are the reasons. But how to phrase that? Maybe the arguments are the items I'm emphasizing or the reasons why the proposal is good.} \\
\\ 
\texttt{Alternatively, maybe the ARGUMENT is a set of items that support my proposal. Hmm, but the example in the problem shows ARGUMENT: \{"..."\} with the content being the actual argument text. Wait, looking back at the problem statement: "ARGUMENT: \{'...'\} Defend your last proposal or argue against the player's proposal." So the ARGUMENT is a set containing a single string, which is the actual argument. So I need to put the argument inside the set as a string.} \\
\\ 
\texttt{So, for example, if I propose a set, my argument could be that the set maximizes my importance while staying under the limit. Let me try that.} \\
\\ 
\texttt{So, my proposal includes A11, B33, B21, B31, C08, B09, A60, A21. Let me check the total effort again. A11:} \\
            }
        }
    }
    & & \\ \\

    \theutterance \stepcounter{utterance}  
    & & & \multicolumn{2}{p{0.3\linewidth}}{
        \cellcolor[rgb]{0.9,0.9,0.9}{
            \makecell[{{p{\linewidth}}}]{
                \texttt{\tiny{[GM$|$GM]}}
                \texttt{abort} \\
            }
        }
    }
    & & \\ \\

    \theutterance \stepcounter{utterance}  
    & & & \multicolumn{2}{p{0.3\linewidth}}{
        \cellcolor[rgb]{0.9,0.9,0.9}{
            \makecell[{{p{\linewidth}}}]{
                \texttt{\tiny{[GM$|$GM]}}
                \texttt{end game} \\
            }
        }
    }
    & & \\ \\

\end{supertabular}
}

\end{document}
